\documentclass[10pt,a4paper]{amsart}
\usepackage[margin=19mm]{geometry}
\usepackage{amsmath,amssymb,mathtools}
\usepackage[scr=rsfs,cal=euler]{mathalfa}
\usepackage{xfrac}
\usepackage[unicode,hidelinks,bookmarks=false]{hyperref}
\newcommand{\ie}{i.e.\ }
\newcommand{\cf}{cf.\ }
\newcommand{\resp}{respectively\ }
\newcommand{\Subsection}{\subsection}
\providecommand{\nobreakdash}{\nobreak\hbox{-}}
\setlength{\emergencystretch}{2em}
\multlinegap0pt
\usepackage{cleveref}

\theoremstyle{plain}
\newtheorem{theorem}{Theorem}[section]
\newtheorem{proposition}[theorem]{Proposition}
\newtheorem{lemma}[theorem]{Lemma}
\newtheorem{corollary}[theorem]{Corollary}
\title[A bandwidth refinement of the Erd\H{o}s distinct subset sums bound]{A bandwidth refinement of the Erd\H{o}s distinct subset sums bound}
\author{Simone Costa, Stefano Della Fiore}
\date{Source: arXiv:2609.27941v1 (CC BY 4.0)}
\begin{document}
\maketitle

\noindent\emph{Source and licence.} The delivered work is the article shown in the title above, version arXiv:2609.27941v1 (CC BY 4.0), distributed under the Creative Commons Attribution 4.0 International licence, as recorded in the arXiv licence metadata for this exact version and shown on the abstract page saved with this item. The full licence notice, which carries the licence URL and the address of that abstract page, is the bundled file \texttt{licenses/arxiv-2609.27941-CC-BY-4.0.txt}.
\par\smallskip

\noindent\textit{Editorial scope.} Counted unit: the article's Lemma (source label lem:Catalan-tail) (source0.tex lines 401--406, source label \texttt{lem:Catalan-tail}): ``As , r ( 43 O(r -1 ) )C r-1 .''. It is delivered together with the article's complete original proof (source0.tex lines 408--509, one proof environment adjacent to the statement, 102 lines), copied line for line from the article's own text. Required context retained uncounted: none beyond the counted statement and proof. Editorial formatting only: an AMS \texttt{amsart} preamble that restates the article's own macro definitions and its own theorem declarations, and the theorem environments the retained lines use; the packages cleveref are loaded to match the article's own setup; the article's own class file is neither shipped nor input; the retained environments are renumbered sequentially in order of appearance, whereas the article prints them with its own numbering; the article's equation numbering is not reproduced and every cross-reference inside the retained text resolves to the wrapper's numbering. No citation occurs inside any retained line, so the wrapper carries no bibliography; All retained bodies are exact copies of the bundled \texttt{sources/211531/source0.tex}, so no omission receipt applies. No asset-inclusion command occurs inside any retained span and the package ships no image asset.


\section{The counted unit: Lemma (source label lem:Catalan-tail) and its complete original proof}

\begin{lemma}\label{lem:Catalan-tail}
As $r\to\infty$,
\begin{equation}\label{eq:Catalan-tail}
 \Gamma_r=\left(\frac43+O(r^{-1})\right)C_{r-1}.
\end{equation}
\end{lemma}

\begin{proof}
For $r\ge2$, define
\[
 q_r:=\frac{\Gamma_r}{C_{r-1}}.
\]
Since
\[
 \Gamma_{r+1}=\Gamma_r+C_r,
\]
we have
\begin{align*}
 q_{r+1}
 &=\frac{\Gamma_{r+1}}{C_r}
 =\frac{\Gamma_r+C_r}{C_r}
 =1+\frac{C_{r-1}}{C_r}\frac{\Gamma_r}{C_{r-1}}=1+\alpha_r q_r,
\end{align*}
where
\begin{align}
 \alpha_r
 :=\frac{C_{r-1}}{C_r}
 =\frac{\frac1r\binom{2r-2}{r-1}}
         {\frac1{r+1}\binom{2r}{r}}
 =\frac{r+1}{r}\frac{r^2}{(2r)(2r-1)}
 =\frac{r+1}{2(2r-1)}.
 \label{eq:alpha-r}
\end{align}
Thus
\begin{equation}\label{eq:q-rec}
 q_{r+1}=1+\alpha_rq_r.
\end{equation}

Write
\[
 q_r=\frac43+e_r.
\]
Substituting this into \cref{eq:q-rec} gives
\begin{align*}
 e_{r+1}
 =1+\alpha_r\left(\frac43+e_r\right)-\frac43
 =\alpha_re_r+\left(\frac{4\alpha_r}{3}-\frac13\right).
\end{align*}
Using \cref{eq:alpha-r},
\begin{align*}
 \frac{4\alpha_r}{3}-\frac13
 =\frac{2(r+1)}{3(2r-1)}-\frac13
 =\frac{2r+2-(2r-1)}{3(2r-1)}
 =\frac{1}{2r-1}.
\end{align*}
Therefore
\begin{equation}\label{eq:e-rec}
 e_{r+1}=\alpha_re_r+\frac{1}{2r-1}.
\end{equation}
For $r\ge2$,
\[
 0<\alpha_r=\frac{r+1}{2(2r-1)}\le\frac12.
\]
We now prove by induction that
\begin{equation}\label{eq:e-bound}
 |e_r|\le\frac4r\qquad(r\ge2).
\end{equation}
For $r=2$,
\[
 \Gamma_2=C_1=1,
 \qquad
 q_2=1,
 \qquad
 e_2=-\frac13,
\]
so \cref{eq:e-bound} holds.  Assume it holds for some $r\ge2$.  Then
\cref{eq:e-rec} and $\alpha_r\le1/2$ give
\[
 |e_{r+1}|
 \le\frac12\frac4r+\frac{1}{2r-1}
 =\frac2r+\frac{1}{2r-1}.
\]
It remains to compare this with $4/(r+1)$.  After subtracting $2/r$ from
both sides, the required inequality is
\[
 \frac{1}{2r-1}\le\frac4{r+1}-\frac2r
 =\frac{2(r-1)}{r(r+1)}.
\]
Multiplying by the positive quantity $r(r+1)(2r-1)$, this is equivalent to
\[
 r(r+1)\le2(r-1)(2r-1).
\]
The difference between the right- and left-hand sides is
\[
 2(r-1)(2r-1)-r(r+1)
 =3r^2-7r+2
 =(3r-1)(r-2),
\]
which is nonnegative for $r\ge2$.  Hence
\[
 |e_{r+1}|\le\frac4{r+1},
\]
and the induction is complete.  Therefore $e_r=O(r^{-1})$, and since
$q_r=4/3+e_r$,
\[
 \Gamma_r=q_rC_{r-1}
 =\left(\frac43+O(r^{-1})\right)C_{r-1}.
\]
\end{proof}

\end{document}
